% Frozen source-aware proof candidate.  Do not insert into the live chapter
% before production Rethlas verification.  No novelty or priority claim is
% intended.

\paragraph{Imported inputs.}
The following primary-source statements are used with their stated
hypotheses.
\begin{enumerate}
\item Gubler--Jell--Rabinoff, Proposition~5.8: if
\(\Lambda=\mathbf Z\), or if \(\Lambda\subset\mathbf R\) is divisible,
then every \(\Lambda\)-PL function is analytically locally a finite
\(\Lambda\)-linear combination of integral PL functions; the converse
holds as well.
\item Gubler--Jell--Rabinoff, \S6.6 and Lemma~6.5: invertible sheaves on
the reduction of a good strict germ are represented on algebraic models,
the models form a filtered system under proper domination, and a finite
family of residue bundles can be represented on one common model.
Proper domination preserves the associated open subset of the
Riemann--Zariski space by the model formalism in their \S6.2.
\item Gubler--Jell--Rabinoff, Lemma~6.18: the residue map on PL germs is
linear.  Their Proposition~7.1 identifies semipositivity of a PL germ at
\(x\) with nefness of its residue class, and their Remark~7.3(4) says that
the semipositivity locus is open.
\item Chambert--Loir--Ducros, Theorem~17.5.7(b), in the version of
25~July~2025: on an almost-good strict germ the strict integral residual
functor is essentially surjective already for analytically trivial
underlying bundles.  In rank one, a residual line bundle is therefore the
residue of an integral PL metric on a locally trivial analytic line bundle.
A good strict germ satisfies the required almost-good hypothesis.  This is
the lifting input used in the proof of Gubler--Jell--Rabinoff,
Proposition~7.10.
\item Chow's lemma, in the form of Stacks Project, Tag~02O2, gives a
proper birational quasi-projective model dominating an integral separated
finite-type model over the residue field.  The ample-twist criterion of
Stacks Project, Tag~0892, says that if \(A\) is ample and \(M\) is an
invertible sheaf on such a model, then \(A^{\otimes N}\otimes M\) is ample
for all sufficiently large \(N\).
\end{enumerate}

\begin{proof}
All assertions are local at \(x\).  Since \(X\) is good and strict, we
may first replace it by a strictly affinoid neighbourhood of \(x\).
This neighbourhood is separated, and its germ at \(x\) is good, hence
almost good in the sense required by
\cite[Theorem~17.5.7(b)]{CLD25}.

Put
\[
                         \mathscr R=\operatorname{red}(X,x).
\]
For each \(i\), the residue
\[
                         L_{h_i}(x)\in\operatorname{Pic}(\mathscr R)
\]
is represented by a line bundle on a model of \(\mathscr R\).  The
description of invertible sheaves in \cite[\S6.6]{GJR26}, together with
the filteredness statement \cite[Lemma~6.5]{GJR26}, therefore gives one
model \(V_0\) of \(\mathscr R\) and line bundles \(M_{i,0}\) on \(V_0\)
representing all the \(L_{h_i}(x)\).

Apply Chow's lemma \cite[Tag~02O2]{stacks-project} to \(V_0\) over
\(\widetilde K\).  It gives a proper surjective morphism
\(Z\to V_0\), with \(Z\) quasi-projective, which is an isomorphism over
a dense open subset of \(V_0\).  Replacing \(Z\) by the reduced closure
of the inverse image of that dense open subset, we obtain an integral
quasi-projective variety \(V\) and a proper birational morphism
\[
                             \pi:V\longrightarrow V_0.
\]
By \cite[\S6.2]{GJR26}, proper domination does not change the associated
open subset of the Riemann--Zariski space.  Hence \(V\) is again a model
of \(\mathscr R\).  Put
\[
                             M_i=\pi^*M_{i,0}.
\]

Choose an ample line bundle \(A\) on \(V\).  Apply the ample-twist criterion
\cite[Tag~0892]{stacks-project} with \(f=\operatorname{id}_V\), with
\(\mathcal M=A\), and successively with
\(\mathcal L=M_i,M_i^{-1}\).  Every invertible sheaf is
\(\operatorname{id}_V\)-ample, since its restriction to an affine open is
ample.  Hence, after taking the maximum of finitely many thresholds, there
is one integer \(N\geq1\) such that
\begin{equation}
\label{eq:p09-ample-order-interval}
                A^{\otimes N}\otimes M_i,
                \qquad
                A^{\otimes N}\otimes M_i^{-1}
\end{equation}
are ample for every \(i\).  In particular, \(A\) and all the line
bundles in \eqref{eq:p09-ample-order-interval} are nef: their restrictions
to every proper integral curve have nonnegative degree.

The line bundle \(A\) determines an invertible sheaf on
\(\mathscr R\).  The strict lifting theorem
\cite[Theorem~17.5.7(b)]{CLD25}, used for line bundles in precisely this
way in the proof of \cite[Proposition~7.10]{GJR26}, gives a strictly
integral PL metric on an analytically trivial line bundle over the germ
\((X,x)\) whose residual line bundle is represented by \((V,A)\).
Choosing a trivializing frame and taking the negative logarithm of its
norm gives an integral PL germ \(a\) satisfying
\[
                             L_a(x)\simeq[A].
\]
After shrinking, regard \(a,h_1,\ldots,h_m\) as actual functions on one
common neighbourhood of \(x\).

Residue linearity \cite[Lemma~6.18]{GJR26} now gives
\[
 \begin{aligned}
 L_{Na+h_i}(x)
   &\simeq [A^{\otimes N}\otimes M_i],\\
 L_{Na-h_i}(x)
   &\simeq [A^{\otimes N}\otimes M_i^{-1}].
 \end{aligned}
\]
These residue classes, as well as \(L_a(x)\), are nef.  The residue
criterion \cite[Proposition~7.1]{GJR26} shows that \(a\) and all
\(Na\pm h_i\) are semipositive at \(x\).  Their semipositivity loci are
open by \cite[Remark~7.3(4)]{GJR26}.  Intersecting these finitely many
open loci and then choosing a strict analytic neighbourhood \(W\) of
\(x\) inside the intersection proves the first assertion.

Let now \(h\) be a \(\Lambda\)-PL germ.  By
\cite[Proposition~5.8]{GJR26}, after shrinking around \(x\), there is an
identity of actual functions
\[
                              h=\sum_{i=1}^{m}\lambda_i h_i,
 \qquad
 \lambda_i\in\Lambda,\quad
 h_i\in\operatorname{PL}_{\mathbf Z}.
\]
Apply the first assertion to \(h_1,\ldots,h_m\), and restrict everything
to the resulting common neighbourhood \(W\).  Omit the indices for
which \(\lambda_i=0\).  For every remaining index, put
\[
 \epsilon_i=
 \begin{cases}
  1,&\lambda_i>0,\\
 -1,&\lambda_i<0,
 \end{cases}
 \qquad
 c_i=|\lambda_i|,
 \qquad
 g_i=Na+\epsilon_i h_i,
\]
and set
\[
                              c=N\sum_i|\lambda_i|.
\]
Every \(g_i\), as well as \(a\), is an integral PL function semipositive
on \(W\).  Moreover,
\[
 \begin{aligned}
 \sum_i c_ig_i-ca
 &=\sum_i|\lambda_i|
       \bigl(Na+\epsilon_i h_i\bigr)
       -N\sum_i|\lambda_i|a\\
 &=\sum_i\lambda_i h_i
 =h.
 \end{aligned}
\]
Since \(\Lambda\) is an additive subgroup of \(\mathbf R\), one has
\[
                 |\lambda_i|\in\Lambda_{\geq0},
 \qquad
                 N\sum_i|\lambda_i|\in\Lambda_{\geq0}.
\]
Thus \eqref{eq:p09-signed-semipositive-span} has all the asserted
coefficient properties.  If all \(\lambda_i\) vanish, take the sum to
be empty and \(c=0\).

Conversely, every finite \(\Lambda\)-linear combination of integral PL
functions is \(\Lambda\)-PL, by the reverse implication in
\cite[Proposition~5.8]{GJR26}.  This proves the final assertion.
\end{proof}

\begin{remark}
\label{rem:p09-signed-semipositive-span-scope}
The proof uses the common residual model only to establish nefness of the
residues of the actual functions \(Na\pm h_i\); it never replaces the
original \(h_i\) by unrelated lifts.  Thus the displayed identity is an
identity of functions, not merely of residue classes.  If \(a'\) is another
integral PL lift of \([A]\), then \(L_{a-a'}(x)=0\) by residue linearity,
and \cite[Lemma~6.18 and Proposition~7.1]{GJR26} identify \(a-a'\) as a
harmonic germ.  In any case, replacing \(a\) by \(a'\) in both occurrences
of \eqref{eq:p09-signed-semipositive-span}, and redefining
\(g_i'=Na'+\epsilon_i h_i\), produces exact cancellation, so there is no
hidden kernel correction.

No interior-point assumption is used.  At a boundary point the first
common residual model need not be proper; Chow's lemma supplies a proper
dominating quasi-projective model, and quasi-projectivity is sufficient for
the ample order-unit argument.

Finally, the theorem is not a positive-cone decomposition theorem.  Its
proof never uses semipositivity of \(h\), and the coefficient \(-c\) in
\eqref{eq:p09-signed-semipositive-span} is generally negative.  It
identifies the signed \(\Lambda\)-span of the semipositive integral germs;
when \(\Lambda=\mathbf Z\), this is their group completion.  It gives no
positive-coefficient decomposition of a semipositive \(\Lambda\)-PL germ.
\end{remark}
